
\subsection{Greedy}

The greedy baseline (we use it as a baseline for comparison the effect of our methods)
corresponds to the \texttt{RLCardAgent} used in our evaluation code.
Unlike the random agent, it uses a small set of hand-crafted rules based on card combinations.
The agent first converts the environment's rank-coded cards into a compact string representation,
where ranks are ordered as

\[
    3,4,5,6,7,8,9,T,J,Q,K,A,2,B,R.
\]

It then distinguishes between two cases: leading a new round and following a
previous move.

When the agent is leading, it decomposes its hand into several structural
combinations using \texttt{combine\_cards}. The decomposition extracts
combinations in a fixed priority order: rocket, bomb, trio, trio chain, solo
chain, pair chain, pair, and finally solo cards.

After this decomposition, the agent searches for a combination containing the
minimum card in its current hand and plays that combination. This produces a
simple greedy leading policy: it tends to release low cards while preserving
larger cards and special combinations for later.

When the agent is following another player's move, it uses RLCard's card-type
table to identify the type and rank of the last move. Among all legal non-pass
actions with the same card type, it chooses the one with the smallest rank that
can beat the previous move. Formally, if \(\mathcal{B}(I)\) is the set of legal
responses with the same type as the last move, the selected action is

\[
    a^* = \arg\min_{a \in \mathcal{B}(I)} \mathrm{rank}(a).
\]

If no such response exists, the agent may pass. In particular, when the previous
valid move was made by the other farmer and the current agent is also a farmer,
it passes to avoid beating its teammate. If none of the above rules applies, the
agent falls back to choosing a random legal action.

This greedy policy is still very limited: it does not perform long-horizon
planning, does not estimate hidden hands, and does not learn from experience.
However, it is stronger than a random baseline because it follows basic
Dou Dizhu principles such as playing the lowest sufficient response, preserving
high cards, and avoiding unnecessary interference with a teammate.
